\documentclass[answers,11pt]{exam}

\usepackage[french]{babel}
\usepackage{ae,lmodern}
\usepackage[utf8]{inputenc}
\usepackage[T1]{fontenc}
\usepackage{amsmath,amssymb,mathrsfs}
\usepackage{graphicx}
\usepackage{enumitem}
\usepackage{stmaryrd}
\usepackage{hyperref}
\usepackage{xcolor}

\newcommand{\ZZ}{\mathbb{Z}}
\newcommand{\KK}{\mathbb{K}}
\newcommand{\FF}{\mathbb{F}}
\newcommand{\RR}{\mathbb{R}}
\newcommand{\CC}{\mathbb{C}}

\title{Leçon 151 : Sous-espaces stables par un endomorphisme ou une famille d'endomorphismes d'un espace vectoriel de dimension finie. Applications}
\author{\sc{Mathieu} Alix}
\date{--}

\begin{document}

\maketitle

\section{Rapport du Jury}

\og Dans cette leçon, il faut présenter des propriétés de l'ensemble des sous-espaces stables par un endomorphisme. Des études détaillées sont les bienvenues, par exemple dans le cas d'une matrice diagonalisable ou dans le cas d'une matrice nilpotente d'indice maximum.\\

L'étude des endomorphismes cycliques et des endomorphismes semi-simples trouvent tout à fait leur place dans cette leçon. Dans le cas des corps $\RR$ ou $\CC$, on pourra, si on le souhaite, caractériser ces derniers par la fermeture de leur orbite. \\

Il ne faut pas oublier d'examiner le cas des sous-espaces stables par des familles d'endomorphismes. Ceci peut déboucher par exemple sur des propriétés des endomorphismes commutant entre eux.\\

La réduction des endomorphismes normaux et l'exemple de résolutions d'équations matricielles peuvent être présentés en applications.\\

La décomposition de Frobenius constitue également une application intéressante de cette leçon. Pour aller plus loin, on peut envisager de développer l'utilisation de sous-espaces stables en théorie des représentations.
\fg

\section{Questions classiques}

Dans tous ce qui suit on considèrera $E$ un $\KK$-ev de dimension finie, sauf si plus de précisions sont données.

\begin{questions}


\question Soit $u \in \mathcal{L}(\RR^3)$ tel que $u^2 \neq 0$ et $u^3=0$ (aka nilpotent d'ordre maximum $3$). Quels sont les sous-espaces stables par $u$?

\begin{solution}
	Soit $e_3$ un vecteur tel que $u^2(e_3) \neq 0$, on pose alors $e_2=u(e_3)$ et $e_1=u^2(e_3)$. La famille $(e_1,e_2,e_3)$ est une base de $E$ (à vérifier). Dans un cette base, $u$ admet la matrice suivante : \[ M=\begin{pmatrix}
		0 & 1 & 0 \\
		0 & 0 & 1 \\
		0 & 0 & 0
	\end{pmatrix} \]
	\begin{itemize}
		\item Si $F$ est un sous-espace stable de dimension $1$, alors il est nécessairement engendré par un vecteur propre, et comme $0$ est la seule valeur propre de $u$ ($u$ nilpotents, alors les vecteurs propres de $F$ sont dans le noyau de $u$. Comme $\ker(u)$ est de dimension $1$ (la matrice $M$ est de rang $2$), alors $F=\ker(u)$.
		\item Si $F$ est sous-espace stable de dimension $2$ alors en considérant l'équation de $F$ : $a_1 x_1 +a_2 x_2 + a_3 x_3=0$. Si $x$ est le vecteur de coordonnées $(x_1,x_2,x_3)$, alors $u(x)$ a pour coordonnées $(x_2,x_3,0)$ et doit vérifier $a_1 x_2 + a_2 x_3=0$ (car $u(F) \subseteq F$) et $a_1 x_3 = 0$ (car $u^2(F) \subseteq F)$. Si $a_1 \neq 0$, alors $x_3=0$ et en reportant dans les $2$ équations précédentes, on trouve $x_2=x_1=0$ ce qui est impossible car dans ce cas $F=\{0\}$. On a donc $a_1=0$, de même $a_2=0$ et donc l'équation de $F$ est $x_3=0$, i.e. $F=\mathrm{Im}(u)$.
	\end{itemize}
	En conclusion les sev stables par $u$ sont $\{0\}
	, \ker(u), \mathrm{Im}(u), \RR^3$.
\end{solution}

\question Soit $u$ un endormorphisme. Montrer que si $\ker(u)$ admet un supplémentaire $S$ stable par $u$, alors $S=\mathrm{Im}(u)$.

\begin{solution}
	Soit $S$ un supplémentaire de $\ker(u)$ stable par $u$. Comme $\ker(u) \cap S = \{0\}$, alors $u$ réalise un isomorphisme de $S$ dans $S$. Ainsi $S=u(S) \subseteq \mathrm{Im}(u)$. Par ailleurs la formule du rang donne \[ \dim(\mathrm{Im}(u))=\dim(E)-\dim(\ker(u))=\dim(S) \]
	D'où l'égalité.
\end{solution}

\question Soit $D$ une droite vectorielle de $\RR^3$. Déterminer les endomorphismes de $\RR^3$ laissant stables tous les plans de $\RR^3$ contenant $D$.

\begin{solution}
	Considérons $e_1$ un vecteur directeur de $D$, qu'on complète en une base $(e_1,e_2,e_3)$ de $\RR^3$. Soit $u$ un endo comme attendu dans l'énoncé. $u$ laisse donc stable $\mathrm{Vect}(e_1,e_2)$ et $\mathrm{Vect}(e_1,e_3)$ et donc laisse stable $D$ qui est leur intersection. La matrice de cet endo est donc de la forme \[ \begin{pmatrix}
		\lambda & \alpha & \gamma \\
		0 & \beta & 0 \\
		0 & 0 & \delta
	\end{pmatrix}
	 \]
	 Enfin $u$ laisse stable le plan $\mathrm{Vect}(e_1,e_2+e_3)$ et, avec mes notations, $u(e_2+e_3)=(\alpha + \gamma)e_1 + \beta e_2 + \delta e_3$ d'où $\beta=\delta$. On vérifie réciproquement qu'un endomorphisme $u$ admettant dans la base $(e_1,e_2,e_3)$ une matrice de la forme \[ \begin{pmatrix}
		\lambda & \alpha & \gamma \\
		0 & \beta & 0 \\
		0 & 0 & \beta
	\end{pmatrix}
	 \]
	 laisse stable tous les plans contenant $D$.
\end{solution}

\question 
\begin{parts}
\part Soit $u \in \mathcal{L}(E)$, montrer que les assertions suivantes sont équivalentes :
\begin{enumerate}[label=\roman*)]
	\item $u$ est une homothétie.
	\item $u$ laisse stable toutes les droites de $E$.
	\item $u$ laisse stable tout hyperplan de $E$.
	\item $u^{\intercal}$ laisse stable toute droite de $E^*$.
	\item $u^{\intercal}$ est une homothétie.
\end{enumerate}
\part On suppose que $E$ est de dimension $n$. que dire d'un $u\in \mathcal{L}(E)$ qui laisse stable tous les sev de dimension $r$ de $E$ pour $r \in \llbracket 1,n-1 \rrbracket$.
\end{parts}

\begin{solution}
	\begin{enumerate}[label=\alph*)]
		\item On effectue l'équivalence des 2 premières, puis un raisonnement circulaire sur les restantes \begin{itemize}
			\item[$\romannumeral 1) \implies \romannumeral 2)$] Si $u$ est une homothétie, il laisse stable tout sous-espace vectoriel de $E$, donc en particulier toute droite de $E$.
			\item[$\romannumeral 2) \implies \romannumeral 1)$] Comme $u$ laisse stable toute droite de $E$, alors pour tout $x \in E$, il existe $\lambda_x \in \KK$ tel que  $u(x)=\lambda_x x$. On souhaite montrer que $\lambda_x$ est indépendant de $x$. Soit $y \in E\backslash\{0\}$,
			\begin{itemize}[label=$\ast$]
				\item Si $\exists \alpha \in \KK, y=\alpha x $, alors $u(y)=\lambda_y y = u(\alpha x)= \alpha u(x)=\lambda_x \alpha x =\lambda_x y $, alors comme $y\neq 0$ alors $\lambda_x=\lambda_y$.
				\item Sinon en regardant $u(x+y)=\lambda_{x+y}(x+y)=u(x)+u(y)=\lambda_x x+\lambda_y y$. Alors $(\lambda_{x+y} - \lambda_x)x = (\lambda_y - \lambda_{x+y})y$, et comme $x$ et $y$ sont linéairement indépendant alors $\lambda_{x+y}=\lambda_x=\lambda_y$.
			\end{itemize}
			Ainsi $\lambda_x$ est indépendant de $x$, et $u$ est une homothétie.
			\item[$\romannumeral 1) \implies \romannumeral 3)$] Si $u$ est une homothétie, il laisse stable tout sous-espace vectoriel de $E$, donc en particulier tout hyperplan de $E$.
			\item[$\romannumeral 3) \implies \romannumeral 4)$] Soit $\varphi$ une forme linéaire non nulle et $H=\ker(\varphi)$ un hyperplan de $E$, et $a \in E$ tel que $\varphi(a) \neq 0$ alors $E=\KK a \oplus H$ (en effet $x\in E$ s'écrit $x=\lambda a + h$ avec $\lambda=\frac{\varphi(x)}{\varphi(a)}$ et $h=x-\lambda a \in H$. Comme $u$ laisse stable $h$
			\begin{align*}
				u^{\intercal}(\varphi)(x)&=\varphi(u(x))=\varphi(\lambda u(a)+u(h))=\lambda \varphi(u(a))=\frac{\varphi(x)}{\varphi(a)} \varphi(u(a)) \\
				&=\mu_{\varphi} \varphi(x) \text{ où } \mu_{\varphi}=\frac{\varphi(u(a))}{\varphi(a)} \in \KK
			\end{align*}
			Ainsi $u^{\intercal}(\varphi)=\mu_{\varphi} \varphi$ et donc $u^{\intercal}$ laisse stable toute droite de $E^{*}$.
			\item[$\romannumeral 4) \implies \romannumeral 5)$] Remplacer $E$ par $E^*$, c'est déjà vu.
			\item[$\romannumeral 5) \implies \romannumeral 1)$] Montrons que si $u^{\intercal}=\lambda Id_{E^*}$ alors $u=\lambda Id_{E}$. Supposons par l'absurde que ce ne soit pas le cas, alors il existe $x\in E\backslash\{0\}$ tel que $y=u(x)-\lambda x \neq 0$. On complète $y$ en une base de $E$ que l'on note $\mathcal{B}$, alors en considérant la forme linéaire $\varphi$ définie par $\varphi(y)=1$ et $\varphi(e)=0$ pour $e \in \mathcal{B}\backslash\{y\}$. On vérifie que $(u^{\intercal}(\varphi)-\lambda \varphi)(x)=\varphi(u(x))-\lambda \varphi(x)=\varphi(y)=1$, ce qui contredit l'homothétie de $u^{\intercal}$.
		\end{itemize}
		\item Montrons que $u$ laisse stable tout hyperplan de $E$, et alors par la question précédente, on aura que $u$ est une homothétie. Si $H$ est un hyperplan de base $(e_i)_{1 \leq i \leq n-1}$, alors le sous-espace de $H$ engendré par $\{e_1,\dots,e_{r-1},e_k\}$ pour $r \leq k \leq n-1$ est laissé stable par $u$, donc $u(e_j) \in H$ pour tout $1 \leq j \leq n-1$, donc $H$ est stable par $u$.
	\end{enumerate}
\end{solution}

\question Stabilité sur les orthogonaux, question d'adjoint, et thm spectrale ? comme exemple d'application.

\question ($E$ un $\RR$-ev) Pour tout endormorphisme $u \in \mathcal{L}(E)$, il existe un sous-espace vectoriel $P$ de dimension $1$ ou $2$ stable par $u$

\begin{solution}
	Si $u$ admet une valeur propre réelle $\lambda$, alors pour tout vecteur propre associé $x\in E\backslash\{0\}$, alors la droite $D=\RR x$ est stable par $u$. 
	
	Sinon $n \geq 2$ et le polynôme minimal $\pi_u$ se décompose dans l’anneau factoriel $\RR[X]$ en produit de facteurs irréductibles de degré $2$ (les valeurs propres de $u$ sont
les racines de $\pi_u$). Ce polynôme $\pi_u$ s’écrit donc $\pi_u(X) =(X^2 +bX+c)Q(X)$ avec $b^2-4c< 0$ et $Q (u) \neq 0$ ($\pi_u$ est le polynôme unitaire de plus petit degré
annulant $u$).

De l’égalité $0=\pi_u(u)= u^2 +b u+c Id  \circ Q(u)$ avec $Q(u)\neq 0$,on déduit alors
que l’endomorphisme $u^2 + bu + c Id$ n’est pas injectif, c’est-à-dire que son noyau n’est pas réduit à $\{0\}$ . Pour tout vecteur $x$ non nul dans ce noyau on vérifie alors que $P = \mathrm{Vect}(x, u(x))$ est un sous-espace vectoriel de dimension $2$ stable par $u$. En effet, $P$ est de dimension $2$ puisque $x$ n’est pas vecteur propre de $u$ et avec $u^2 (x)+bu(x)+cx = 0$ on déduit que $u^2(x)$ est dans $P$, ce qui entraîne que $P$ est stable par u.

\end{solution}


\end{questions}

\section{Question plus exotique ?}

\begin{questions}

\question L'objectif de cet exercice est de démontrer le résultat suivant : si $u \in \mathcal{L}(E)$ vérifie $\pi_u$ irréductible alors $u$ est un endomorphisme semi-simple (\emph{i.e.}  tout sev $F$ de E stable par $u$ admet un supplémentaire dans $E$ stable par $u$). Considérons dès lors $F$ un sev stable de $E$ par $u$
 \begin{parts}
 	\part Traiter le cas $F=E$
 	\part Si $F\neq E$, il existe $x_1 \in E\backslash F$, considérons l'ensemble suivant ${E_{x_1,u} = \{P(u)(x_1), P\in \KK[X]\}}$, vérifier qu'il est stable par $u$ ?
 		\begin{enumerate}[label=\roman*)]
 			\item Notons $I_{x_1}=\{P \in \KK[X], P(u)(x_1)=0\}$, Montrer que $I_{x_1}$ est un idéal principal non-nul. On notera pour la suite $\pi_{x_1}$ un générateur.
 			\item Montrer que $\pi_u=\pi_{x_1}$ et en déduire l'irréductibilité de $\pi_{x_1}$.
 			\item Montrer que $E_{x_1,u} \cap F = \{0\} $
 			\item Conclure en fonction de si $E_{x_1,u}$ est supplémentaire par rapport à $F$. Sinon en déduire un procédé qui permet de conclure.
 		\end{enumerate}

 \end{parts}
 \underline{Remarque:} Ce résultat est essentielle pour pouvoir montrer le résultat suivant sur les endomorphismes semi-simples (voir Gourdon, Algèbre et probabilités pour une démonstration) : \og $u$ est semi-simple si et seulement si $\pi_u = M_1 \dots M_r$ où les $M_i$ sont unitaires, irréductibles et distincts \fg
 
 \begin{solution}
 	    $\triangleright$ Soit $F$ un sev stable par $u$. Montrons l'existence d'un supplémentaire $S$ stable par $u$.\\
\begin{enumerate}[label=(\alph*)]
    \item  Si $F=E$, c'est terminé en prenant $S=\{0\}$ qui est bien stable par $u$.
    
    \item Sinon, considérons $x_1 \in E\backslash F$ et posons l'ensemble suivant qui est stable par $u$  \[ E_{x_1,u} = \{P(u)(x_1), P \in \KK[X]\} \]

    Montrons que cet ensemble pourrait être un potentiel candidat de supplémentaire en commençant par montrer le caractère direct. Montrons que $E_{x_1,u} \cap F = \{0\}$. Pour se faire, on s'intéresse à l'ensemble \[I_{x_1} = \{P\in \KK[X], P(u)(x_1)=0\}\] 
    \begin{enumerate}[label=\roman*)]
    \item Il est laissé au lecteur de vérifier que $I_{x_1}$ est un idéal de $\KK[X]$. Or $\KK[X]$ est principal car $\KK$ est un corps, $I_{x_1}$ est donc principal, celui-ci est non réduit à $\{0\}$ car $\pi_u \in I_{x_1}$. Il est donc engendré par un élément $\pi_{x_1} \in \KK[X]$. 
    \item Or comme dit précédemment, on a $\pi_u \in I_{x_1} = (\pi_{x_1})$  donc $\pi_{x_1} \mid \pi_u$ et comme $\pi_u$ est irréductible (et que $\pi_{x_1} \neq 1$ car $x_1 \neq 0 $ étant dans $E\backslash F$) alors $\pi_u = \pi_{x_1}$ et donc $\pi_{x_1}$ est irréductible.
    \item Considérons $y\in E_{x_1,u} \cap F$, alors, il existe $P \in \KK[X], y =P(u)(x_1)$. Maintenant supposons que $y\neq 0$. alors $P\not \in I_{x_1} =(\pi_{x_1})$ donc $\pi_{x_1}$ ne divise pas $P$ et comme $\pi_{x_1}$ est irréductible, alors $\pi_{x_1}$ et $P$ sont premiers entre eux. Comme $\KK[X]$ est principal, par le théorème de Bezout :

    Il existe $U,V \in \KK[X], UP + V\pi_{x_1} = 1$ donc : \[ x_1 = U(u) \circ P(u)(x_1) + V(u) \circ \pi_{x_1}(u)(x_1) = U(u)(y)\]

    Or $y \in F$ et comme $F$ est stable par $u$ alors par l'égalité au dessus, on devrait avoir $x_1 \in F$, ce qui est impossible. d'où le caractère direct.
    \item Si $F\oplus E_{x_1} = E$, c'est terminé en choisissant $S=E_{x_1}$. Sinon, on choisit $x_2 \in E\backslash (E_{x_1} \oplus F)$ et on recommence. On itère ainsi ce procédé qui se termine en un nombre fini d'itération car $E$ est de dimension fini (et que $\dim(E_{x_i}) \geq 1$). Il existe donc $x_1,\dots x_r$ tel que \[ E = F \oplus E_{x_1} \oplus \cdots \oplus E_{x_r} \]

    et comme pour tout $i \in \{1, \cdots, r\}, E_{x_i}$ est stable par $u$ alors $S=E_{x_1} \oplus \cdots \oplus E_{x_r}$ est stable par $u$ et $F\oplus S=E$.
    \end{enumerate} 
   \end{enumerate}
 \end{solution}
 
 
\question (Décomposition de Fitting) : On considère $u$ un endomorphisme sur une espace vectoriel $E$ de dimension finie $d$.
	\begin{parts}
		\part Les suites $(\ker(u^k))_{k\in \mathbb{N}}$ et $(\mathrm{Im}(u^k))_{k\in\mathbb{N}}$ sont respectivement croissantes et décroissantes.
		\part Montrer que ces suites sont stationnaire à partir d'un même rang $n_0$.
		\part Avec les notations précédentes, montrer que $E=\ker(u^{n_0}) \oplus \mathrm{Im}(u^{n_0})$.
		\part Montrer que $u$ laisse stable $\ker(u^{n_0})$ et $\mathrm{Im}(u^{n_0})$.
		\part Montrer que $u$ induit un endomorphisme nilpotent sur $\ker(u^{n_0})$.
		\part Montrer que $u$ induit un automorphisme sur $\mathrm{Im}(u^{n_0})$.
	\end{parts}
\end{questions}

\begin{solution}
	Voir tout libre d'algèbre linéaire, [MAN] ou [CAL - Carnet de voyage en algébrie] ou si vous le souhaitez tout bien détaillé \href{https://agreg-maths.fr/uploads/versions/5366/Dénombrement%20des%20endomorphismes%20nilpotents%20sur%20un%20corps%20fini.pdf}{ici}
\end{solution}




\end{document}

